\subsection{Five and six objects with sixteen arrows}

\begin{theorem}\label{C:sixteen-five-six-objects}
No finite category with five or six objects and at most sixteen
arrows has classifying space homotopy equivalent to
$\mathbb {CP}^2$.
\end{theorem}
\begin{proof}
Theorem~\ref{C:CP2-lower16} treats at most fifteen arrows.
We may assume that the category is connected and has exactly
sixteen arrows.  The following enumeration uses complete
ordinary corner-monoid catalogs, including both orientations;
it does not identify a monoid with its opposite when that
would omit an actual corner choice.

First enumerate the underlying simple undirected support:
an edge records that at least one of the two cross hom-sets
is nonempty.  For $k$ objects this is a connected graph with
at most $16-k$ edges.  Enumerate every graph on $k$ labeled
vertices, and retain one representative of each orbit under
all $k!$ vertex permutations.  The exact inventories are
\[
\begin{array}{c|r|r|r}
 k&\text{labeled supports}&\text{support orbits}&
 \text{raw cardinality decorations}\\ \hline
 5&728&21&9\,039\,898\\
 6&24\,763&94&7\,492\,422
\end{array}
\]
The raw decoration count is before any structural or symmetry
prune.  To explain its completeness, each vertex receives a
positive corner size, and an undirected edge receives an
ordered pair $(a,b)$ of nonnegative cross cardinalities with
$a+b>0$.  Subject only to total arrow count sixteen, all such
choices are tried.  For a fixed support with $e$ edges their
number is the coefficient of $x^{16-k-e}$ in
\[
 (2-x)^e(1-x)^{-k-2e}.
\]
Indeed, a corner's excess contributes $(1-x)^{-1}$ and an
edge's excess contributes
$\sum_{j\geq0}(j+2)x^j=(2-x)/(1-x)^2$.
The full automorphism group of the support, together with
opposition, removes duplicate decorations.  No additional
vertex permutation is needed: any isomorphism of decorated
categories must already be an automorphism of the underlying
support representative.

Apply the previously proved initial/terminal, retract,
singleton cofinal-deletion, and singleton-core prunes.
A five-object corner of order eight forces the remaining
cross support to be a singleton tree and is excluded by
Lemma~\ref{C:tree-singleton-core-cups}; a larger corner,
or an order-eight corner with six objects, cannot fit the
connected arrow budget.  Thus every remaining corner has
order at most seven.  Acyclic choices with semisimple
rational diagonal are excluded by
Lemma~\ref{C:sixteen-high-semisimple}.  The program uses a
nonzero regular-trace determinant modulo $32003$ only as a
sufficient certificate for rational semisimplicity; a zero
determinant never excludes a corner.

For every remaining decoration, insert every literal corner
multiplication table from the complete ordinary catalogs.
Every cross product is then a variable whose possible values
are exactly its target hom-set.  Associativity propagation
assigns a variable only when an already known associativity
equation forces it, and rejects only contradictory equations.
The depth-first search tries every remaining allowed value.
Consequently every category in the retained scope appears.
Permutations inside cross hom-sets remove duplicate finished
tables; omitting corner permutations at this stage can only
retain additional representatives.  An exact normalized-bar
calculation first retains the tables with mod-two homology
$H_1=0$, $H_2=\mathbb F_2$.

The completed searches and their independent literal bar gates
have the following counts:
\[
\begin{array}{c|r|r|r|r}
 k&\text{finished corner-representative tables}&
 \text{retained rows}&u^2=0&H_3(\mathbb F_2)\ne0\\ \hline
 5&8\,301\,658&411\,909&411\,734&175\\
 6&51\,458&4\,022&4\,022&0
\end{array}
\]
For every retained row the bar gate checks all typed category
axioms, computes a nonzero degree-two cocycle, and either
exhibits a degree-three primitive for its square or computes
positive third mod-two homology.  Neither alternative is
compatible with $\mathbb {CP}^2$.  This proves the theorem.
\end{proof}

\begin{remark}[Scope and independent verification]
\label{C:sixteen-five-six-verification}
The producer is \path{c_categories16_support_high.cpp}.
The completed source stems are
\path{c_categories_16_5_support_CP2} and
\path{c_categories_16_6_support_CP2}.
The independent program \path{d_sixteen_support_graph_audit.py}
constructs graph orbits by expanding every orbit, rather than
using the producer's minimum-representative test.  It verifies
all $21$ and $94$ support representatives, their full
automorphism groups, every formal decoration count, and the
complete start/finish support logs.  The earlier
\path{d_sixteen_generator_scope_audit.py} records the
structural-prune and code-difference audit.

The literal normalized-bar verifier
\path{c_category_tsv16_cups.cpp} imports neither the category
generator nor its homology filter.  Its mathematical routines
are the previously independently audited fifteen-arrow
verifier; the sixteen-arrow specialization changes only the
two input-size guards.  Its completed output stems are
\path{c_category_tsv_16_5_support_complete} and
\path{c_category_tsv_16_6_support_complete}.
It checks respectively $124\,599\,941$ and $633\,948$ typed
associativity equations, $63\,327\,254$ and $227\,919$
cocycle equations, and $298\,082\,532$ and $675\,597$
square-primitive equations.

The separate checker \path{d_verify_sixteen_high_linkage.py}
compares every literal source, target, and multiplication table
with its certificate, requires complete source and certificate
EOFs and final producer counts, and hashes the entire source.
It additionally reconstructs the category and normalized bars
from scratch in Python for all $4\,022$ six-object rows and
all $175$ five-object positive-$H_3$ rows, recomputing the
ranks, the nonzero cocycle, and every applicable primitive.
This additional Python check is not claimed for the remaining
five-object rows: their complete independent check is the
C++ bar verifier just described.  The combined exact-linkage
record is \path{d_sixteen_high_linkage_verified.json}.
No partial sixteen-arrow three- or four-object search is used
in this theorem.
\end{remark}
